% ==========================================================
% 090_design_principles_and_nongoals.tex
% Forecast Admissibility Surface (FAS)
% ==========================================================

\section{Design Principles and Non-Goals}
\label{sec:fas_design_nongoals}

The Forecast Admissibility Surface is intentionally narrow in scope. Its design
reflects a deliberate separation between structural governance decisions and
downstream modeling, evaluation, and optimization. This section articulates the
core principles that guided the design of FAS and clarifies what FAS is not
intended to do.

\subsection{Design Principles}

\paragraph{Determinism.}
FAS produces the same admissibility surface whenever it is applied to the same
inputs. There are no learned parameters, stochastic procedures, or adaptive
mechanisms. This determinism is essential for reproducibility, governance, and
auditability.

\paragraph{Baseline-Centricity.}
FAS relies exclusively on baseline-derived error anatomy. This ensures that
admissibility decisions reflect structural properties of the slice rather than
artifacts of model sophistication. Baseline models are used as probes, not as
performance benchmarks.

\paragraph{Slice-Local Decisions.}
Each slice is evaluated independently. Admissibility is not comparative across
slices, nor is it influenced by global rankings or resource allocation
considerations. This preserves clarity of interpretation and avoids implicit
prioritization.

\paragraph{Conservatism Under Uncertainty.}
When evidence is insufficient to support a definitive classification, FAS errs
toward \Conditional\ rather than \Blocked. Lack of data is treated as uncertainty,
not hostility.

\paragraph{Joinability and Transparency.}
Admissibility outcomes are produced at explicit slice keys and are designed to be
joined deterministically onto downstream artifacts. All inputs and thresholds
used to produce the surface are observable and inspectable.

\subsection{Non-Goals}

\paragraph{Not a Performance Metric.}
FAS does not measure forecast accuracy, cost efficiency, or service quality. It
does not rank slices or models and should not be interpreted as a proxy for
forecast value.

\paragraph{Not a Model Selector.}
FAS does not choose between forecasting models, tune hyperparameters, or guide
model architecture decisions. Such choices belong to downstream modeling and
evaluation processes.

\paragraph{Not an Optimization Objective.}
FAS is not optimized, calibrated, or learned from data. There is no loss
function, gradient, or target distribution associated with admissibility.

\paragraph{Not a Substitute for Governance.}
While FAS is a governance primitive, it does not replace broader governance
policies, human judgment, or operational decision-making. Blocking a slice from
forecasting does not resolve the underlying operational challenge; it merely
prevents unstable forecasts from compounding it.

By adhering to these principles and boundaries, FAS remains a focused and
defensible component of the Forecast Readiness Framework, providing structural
clarity without encroaching on downstream responsibilities.
